\cotefile{34}
\scanpage{116}
\mspage{31}
% archive sheet number: 115

% L01-04
\uline{Corollaire} 4.\unsure{7}. \unsure{Dans|Sous} les conditions de 4.4, supposons $X\inserted{=\operatorname{Spec} A}$ \unsure{strictement} local et satisfaisant au\deleted{x conditions du} th. de pureté cohomologique. Soit $m$ un \deleted{\unsure{entier, qui}} multiple de $n$, premier à la car. résiduelle de $X$, et soit

% L05
\[
\unsure{X'_m|X'}=\operatorname{Spec} A[T_1,\ldots,T_r]/(T_1^m-\unsure{a_1|g_1},\ldots,T_r^m-a_r),
\]

% L06-07
où \deleted{\unsure{pour}} les $a_i$ sont des équations \deleted{\unsure{les}} des $D_i$ $(1\leq i\leq r)$. Soit $U'=X'\times_X U$, Alors l'homomorphisme \unsure{canonique}

% L08
\[
H^*(U,\Lambda)\unsure{{-}|\longrightarrow}H^*(U',\Lambda)
\]

% L09
est nul en dimensions $>0$.

% L10-12
En effet, c'est trivial en dim. 1, car (4.5.1) \deleted{\unsure{dit}}\inserted{équivaut} précisément \inserted{à dire} \deleted{\unsure{que nous}} que $U'$ est un revêtement \deleted{\unsure{galois}} principal connexe, de groupe $\unsure{\mu_n|\mu_m}$.

% L13-14
\deleted{4.8. Soit maintenant $X$ \unsure{str.} local régulier satisfaisant au th. de pureté cohomologique.}

% L15-18
4.8. Soit $X\inserted{=\operatorname{Spec} A}$ un schéma strictement local régulier, $(a_i)_{1\leq i\leq r}$ une partie d'un système régulier de paramètres de $A$, de sorte que \deleted{\unsure{les}} \deleted{$D=\operatorname{div}(a_i)$}

% L18
\[
\boxed{D=\operatorname{div}\left(\prod a_i\right)=\sum D_i\qquad (\text{où }D_i=\operatorname{div}a_i)}\quad\deleted{\unsure{est}}
\]

% L19-20
est un diviseur à croisements normaux, à composantes premières $D_i$ régulières. Posons encore

% L21
\[
\boxed{U=X-\operatorname{supp}D=X-\bigcup D_i},
\]

% L22-24
\deleted{\illegible{1}} et pour tout entier $m$ premier à la car. résiduelle de $X$, introduisons
\[
X(m)=\operatorname{Spec}A[T_1,\ldots,T_r]/(T_1^m-a_1,\ldots,T_r^m-a_r),
\]
